%
% This Section is written in LaTeX because Markdown wasn't playing well with the tables
%
\newpage
%%%%%%%%%%%%%%%%%%%%%%%%%%%%%%%%%%%%%%%%%%%%%%%%%%%%%%
\section{Tables}
\label{sec:tables}
%%%%%%%%%%%%%%%%%%%%%%%%%%%%%%%%%%%%%%%%%%%%%%%%%%%%%%

As a rule of thumb, if you can show your data with a table, consider doing that first. Tables are quite effective at `showing your data', and are usually easy to interpret. The Vis literature even suggests that tables might be more convincing for a skeptical audience \citep{pandey2014persuasive}!

In terms of table formatting, many of the same principles apply to tables as figures. My advice is to: 

\begin{enumerate}
    \item Maintain a high Data-Ink Ratio \citep{tufte1983visual}: 
        \begin{itemize}
            \item Don't use vertical lines
            \item Use horizontal lines sparingly
            \item Use whitespace to help group things together that belong together
        \end{itemize}
    \item Use \texttt{\textbackslash toprule}, \texttt{\textbackslash midrule}, and \texttt{\textbackslash bottomrule} to frame your table
    \item Use the \texttt{[tb!]} flags when declaring the table in \LaTeX to push it to the top or bottom of a page
    \item Use the \texttt{label\{\}} command to give your table a label that you can use later with \texttt{autoref\{\}}
\end{enumerate}
%
In general, writing tables directly in \LaTeX is a huge pain --- the output is generally \emph{beautiful}, but the input is complicated, and generally not very human-readable. I suggest using an online table generator, like \href{https://www.tablesgenerator.com}{https://www.tablesgenerator.com}, and then fine-tuning the formatting in your document. \autoref{tab:exampletable} is an example from \citet{wallace2013}. 

\begin{center}
\begin{table*}[b!]
\small
\begin{center}
\begin{tabular}{l l l l l}
\toprule
Condition & \multicolumn{4}{l}{Teamwork Measures}\\
\cmidrule(r){2-5}
& \multicolumn{2}{l}{Measure 1} & Measure 2 & Measure 3 \\
\cmidrule(r){2-3} 
&  Sub 1 & Sub 2 &  \\ 


\midrule


Condition 1 & .478 & .468 & 1.68 & 12.3  \\
                  & (.0259) & (.031) & (.703) & (5.93) \\
\\

Condition 2 & .444 & .451 & 1.32 & 18.4 \\
                              & (.0478) & (.0585) & (.319) & (2.76) \\
\\

Condition 3 &  & .410 &  2.857 &  \\
                     &        & (.0384)  & (.350) & \\
\\


ANOVA Results & $F_{(1,12)} = 2.27,$ & $F_{(2,18)} = 2.76,$  & $F_{(2,18)} = 16.16,$ & $F_{(1,12)} = 6.16,$ \\
                          & $p = .158$ & $p = .090  $ &  $p < .0001$* &  $p = 0.029$ * \\
\bottomrule
\end{tabular}

\caption[a]{Mean values and standard deviations (in parentheses) for teamwork measures, and ANOVA results for comparisons between experimental conditions. Significant results denoted by *.}
\label{tab:exampletable}
\end{center}
\end{table*}
\end{center}

\subsection*{Additional Resources}
\begin{enumerate}
    \item Overleaf's documentation for tables: \href{https://www.overleaf.com/learn/latex/Tables}{https://www.overleaf.com/learn/latex/Tables}
\end{enumerate}


%
% ADAPT 
%


% \definecolor{palered}{HTML}{ffadad}
% \definecolor{paleorange}{HTML}{ffd6a5}
% \definecolor{paleyellow}{HTML}{fdffb6}
% \definecolor{palegreen}{HTML}{caffbf}
% \definecolor{palecyan}{HTML}{9bf6ff}
% \definecolor{paleblue}{HTML}{a0c4ff}
% \definecolor{palepurple}{HTML}{bdb2ff}
% \definecolor{palepink}{HTML}{ffc6ff}
% \definecolor{palewhite}{HTML}{fffffc}

% \begin{table}[bt!]
%     \centering
%     \small
    
% \begin{tabular}{p{4.5cm}p{3cm}p{3cm}p{3cm}}
% \toprule
% \textbf{Main Themes} & \textbf{Sub Themes} \\ 
% \midrule
% \rowcolor{palered!10}\textcolor{palered}{\faCircle}~\textbf{(1)} Decisions to play are based on childhood experiences relative to the cultural norm growing up  &  \textcolor{palered}{\faChevronCircleRight}~\emph{Nostalgic memories  and tabletop games}  & \textcolor{palered}{\faChevronCircleRight}~\emph{Families with low SES} & \\ 

% \rowcolor{paleorange!10}\textcolor{paleorange}{\faCircle}~\textbf{(2)} Adapted and influenced cultural themes in tabletop games should teach more about culture &  \textcolor{paleorange}{\faChevronCircleRight}~\emph{Participants call on game designers} & & \\ 

% \rowcolor{paleblue!10}\textcolor{paleblue}{\faCircle}~\textbf{(3)} Having someone in mind to play with is the strongest contribution to purchase decision of a tabletop game &  \textcolor{paleblue}{\faChevronCircleRight}~\emph{Interested players and limited players}  & \textcolor{paleblue}{\faChevronCircleRight}~\emph{Word of mouth through price, reviews, and demonstration} & \textcolor{paleblue}{\faChevronCircleRight}~\emph{Quality of game components} \\ 

% \rowcolor{palegreen!10}\textcolor{palegreen}{\faCircle}~\textbf{(4)} Bonding experiences emerge through collaborative and competitive playstyle regardless of household income &  \textcolor{palegreen}{\faChevronCircleRight}~\emph{Disengaged from the digital world} & & \\ 

% \rowcolor{palecyan!10}\textcolor{palecyan}{\faCircle}~\textbf{(5)} Portability is connected for playing  frequency of a tabletop game &  \textcolor{palecyan}{\faChevronCircleRight}~\emph{Playing space and size matters} & \textcolor{palecyan}{\faChevronCircleRight}~\emph{Storage space} &  \textcolor{palecyan}{\faChevronCircleRight}~\emph{High-SES families prefer games with a higher learning curve} \\ 
% \rowcolor{palepurple!10}\textcolor{palepurple}{\faCircle}~\textbf{(6)} Tabletop games should have app-assisted technology   & \textcolor{palepurple}{\faChevronCircleRight}~\emph{Tutorial that demonstrates how the game should be played} &  \textcolor{palepurple}{\faChevronCircleRight}~\emph{Game rules are confusing} &  \\ 

% \rowcolor{palepink!10}\textcolor{palepink}{\faCircle}~\textbf{(7)} Involvement with tabletop games community culture  &  \textcolor{palepink}{\faChevronCircleRight}~\emph{Community involvement} & & \\ 
% \bottomrule
% \end{tabular}

% \caption{We inferred seven themes from our hybrid (largely reflexive) thematic analysis (drawing from codebook approaches) that are relevant to cultural influences, play and purchase decisions in tabletop games. We present a complete overview of the themes and extrapolated subthemes in this table.}
% \label{tab:findings}
% \end{table}
